%==============================================================================
%  03-value-ladder.tex — Section 2: The AI Value Creation Ladder
%==============================================================================
\strategyPart{07 \textbar\ The AI Value Creation Ladder}{The AI Value Creation Ladder}

Value from AI is not binary. It accumulates in distinct levels, and each level
places different demands on the organization. Gartner's AI maturity research
describes the organizational progression — from \textit{Aware} to
\textit{Active}, \textit{Operational}, \textit{Systemic} and
\textit{Transformational} [Gartner AI Maturity Model] — and this section
translates that progression into a concrete value ladder that CEOs can use to
diagnose where their organization creates value today.

\section{The four levels of value creation}

\begin{center}
  \renewcommand{\arraystretch}{1.5}
  \rowcolors{2}{inSnow}{white}
  \fitwidth{\begin{tabular}{>{\raggedright\arraybackslash}p{2.1cm}>{\raggedright\arraybackslash}p{3.2cm}>{\raggedright\arraybackslash}p{6.3cm}}
    \rowcolor{inNavy}
    \hcell{Level} & \hcell{Mode} & \hcell{What the organization gains} \\
    \textbf{1} & \textbf{Assist} — augmentation & Human expertise is amplified: drafting, search, analysis, recommendations. Humans stay in control. \\
    \textbf{2} & \textbf{Automate} — process efficiency & Structured workflows run without human touch: claims, onboarding, forecasting, routing. \\
    \textbf{3} & \textbf{Agentic} — autonomous execution & AI agents plan and execute multi-step tasks with defined guardrails and human oversight. \\
    \textbf{4} & \textbf{Redesign} — business model reinvention & Products, revenue models and value propositions are rebuilt around AI-native capabilities. \\
  \end{tabular}}
\end{center}

\vspace{0.6em}
\begin{center}
  \fitwidth{\begin{tikzpicture}
    \node[fill=inLight,  text=white, font=\bfseries\footnotesize, minimum width=3.2cm, minimum height=1.0cm, align=center] (l1) at (0,0) {L1 \\ Assist};
    \node[fill=inBright, text=white, font=\bfseries\footnotesize, minimum width=3.2cm, minimum height=1.4cm, align=center] (l2) at (0,1.2) {L2 \\ Automate};
    \node[fill=inHover,  text=white, font=\bfseries\footnotesize, minimum width=3.2cm, minimum height=1.8cm, align=center] (l3) at (0,2.8) {L3 \\ Agentic};
    \node[fill=inNavy,   text=white, font=\bfseries\footnotesize, minimum width=3.2cm, minimum height=2.2cm, align=center] (l4) at (0,4.8) {L4 \\ Redesign};
    \node[font=\scriptsize\bfseries, color=inGrey, anchor=west] at (1.8,0.5) {Augmentation — humans + AI};
    \node[font=\scriptsize\bfseries, color=inGrey, anchor=west] at (1.8,1.9) {Process efficiency — workflow automation};
    \node[font=\scriptsize\bfseries, color=inGrey, anchor=west] at (1.8,3.7) {Autonomous execution — agents};
    \node[font=\scriptsize\bfseries, color=inGrey, anchor=west] at (1.8,5.9) {Business-model reinvention};
  \end{tikzpicture}}
\end{center}

\section{Why the levels must be sequenced}

Each level assumes the one below it. Level 3 (agentic) without a Level 1
foundation of trust and a Level 2 foundation of reliable data produces
autonomous failures at scale. The research quantifies how quickly this frontier
is approaching: Gartner predicts that by 2028, 33\% of enterprise software
applications will include agentic AI — up from less than 1\% in 2024 —
enabling 15\% of day-to-day work decisions to be made autonomously
[Gartner 2024].

\section{The diagnostic question for the executive}

For every business process, ask: \textbf{where on the ladder do we operate
today, and what is the next rung?} Most organizations discover that their AI
portfolio is crowded at Level 1 with isolated proofs, empty at Level 2 where
processes should be redesigned, and ungoverned at the frontier of Levels 3--4.
The Value Creation Ladder converts that discovery into a sequenced investment
agenda rather than a scramble of experiments.

\takeaway{Sequence value creation: assist, automate, agentic, redesign — each level builds on the foundations below it.}
